\MakeAbstract{
针对电缆通道空间狭窄、障碍物较多的特点，本文选用五连杆轮腿机器人作为检测平台，并分别研究其运动控制和局部路径规划方法。建模时将左右腿和机体偏航运动统一纳入状态空间，再把单侧五连杆机构等效为虚拟腿。控制器由LQR与VMC串联构成：LQR计算车体所需的控制量，VMC负责将其转换为轮毂及关节电机的输出。MATLAB/Simulink仿真中，机器人连续越过了四个高度为0.25~m的障碍物，越障过程中机体姿态保持稳定。

路径规划实验在Sparrow环境中完成。本文先以DWA作为传统方法基线，再训练DQN和SAC，并围绕动作空间、奖励函数及训练参数进行调试。为减轻环境采样与网络更新相互等待的问题，在SAC训练流程中接入Actor--Learner--Sharer架构。随后将连续多帧激光雷达观测送入Transformer编码器，并与SAC组合为TransSAC，使策略能够利用障碍物的历史运动信息。在相同训练条件下，TransSAC较早进入稳定阶段，训练末期到达率为0.87，任务步数也低于其余对比算法。
}{运动控制，深度强化学习，路径规划，TransSAC}
    
\MakeAbstractEng{
This graduation design investigates motion control and path planning for a wheel-legged cable-channel inspection robot. A unified dynamic model incorporates both legs and the yaw motion to reduce the coordination problems associated with independent leg control. Each five-link leg is represented as a virtual leg, and an LQR controller combined with VMC maps the desired body motion to the wheel and joint motor inputs. MATLAB/Simulink simulations show that the robot can continuously traverse four 0.25~m obstacles while maintaining a stable posture.

For path planning, DWA, DQN, and SAC are evaluated in the open-source Sparrow environment. An Actor--Learner--Sharer architecture is introduced to improve data throughput during SAC training. TransSAC further combines a Transformer encoder with SAC to incorporate historical LiDAR observations. Under the same training conditions, TransSAC converges faster, reaches a final arrival rate of 0.87, and achieves better path efficiency.
}{motion control, deep reinforcement learning, path planning, TransSAC}
